% ---------------------------------------------------------------------
%  Chapitre 4 — Le contexte des trades : comparer et caractériser
%  (ex-étapes 2 à 6 du Bloc 1 : ANOVA, khi-deux, ACP, AFC, ACM)
%  Rôle dans le mémoire : une fois le jugement individuel rendu
%  (chap. 3), on compare les stratégies entre elles et on cartographie
%  les situations d'entrée. Analyses DESCRIPTIVES et comparatives —
%  les réserves d'indépendance du chap. 3 s'appliquent et sont
%  signalées dans les encadrés.
% ---------------------------------------------------------------------
\chapter{Le contexte des trades~: comparer et caractériser}
\label{chap:contexte}

% [À RÉDIGER — chapeau du chapitre, 3-4 lignes] Idées :
%  - Le chap. 3 a jugé chaque stratégie isolément (contre le hasard).
%    Ce chapitre change de question : les stratégies diffèrent-elles
%    ENTRE ELLES ? et le succès dépend-il du CONTEXTE (régime, secteur,
%    situation d'entrée) ?
%  - Statut des conclusions : comparatif et descriptif — on n'y valide
%    aucune stratégie (le verdict appartient au chap. 3).

\begin{limites}
% [À RÉDIGER quand on traitera ce chapitre — pour l'instant on ÉCARTE l'oracle]
La stratégie \emph{oracle} (témoin, chapitre~\ref{chap:juger}) est
\textbf{exclue de ce chapitre}~: son avantage artificiel
(\edgevar{} moyen $\approx +26\,\%$) écraserait toute comparaison
(analyse de la variance, khi-deux, plans factoriels). Elle n'a de sens
que comme témoin de \emph{validation}, non comme stratégie à comparer aux
autres. Les analyses ci-dessous portent donc sur les \textbf{dix} vraies
stratégies. % chiffres des étapes 2-6 à régénérer sur ces 10 (voir TODO)
\end{limites}

% =====================================================================
\section{Les stratégies diffèrent-elles entre elles ?}
\label{sec:b1-etape2}

\subsection{Présentation théorique}
% [À DÉVELOPPER] ANOVA à 1 facteur : décomposition SCE_inter / SCE_intra,
% statistique F ~ Fisher(k-1, n-k). Puis Tukey HSD (range studentisé).
% Puis ANOVA à 2 facteurs (strategy x regime) avec terme d'interaction.
\paragraph{ANOVA à un facteur.}
\[
  F = \frac{\mathrm{SCE}_{\text{inter}}/(k-1)}
           {\mathrm{SCE}_{\text{intra}}/(n-k)} \sim F_{k-1,\,n-k}.
\]
La décomposition de la variance et le test \(F\) suivent le cadre
classique de l'analyse de la variance~\cite{scheffe1959}.

\paragraph{Comparaisons multiples (Tukey HSD).}
Fondées sur la loi du \emph{range studentisé}, elles corrigent
automatiquement le risque lié aux $k(k-1)/2$ paires.

\paragraph{ANOVA à deux facteurs avec interaction.}
Modèle \(\edgevar = \mu + \alpha_i + \beta_j + \gamma_{ij} + \varepsilon\)~;
le terme d'interaction \(\gamma_{ij}\) teste si l'effet de la stratégie
\emph{dépend du régime de marché}.

\begin{conditions}
Indépendance des observations --- hypothèse mise à l'épreuve au
chapitre~\ref{chap:juger}~: les p-values de ce chapitre, calculées au
niveau trade, s'entendent donc sous cette réserve (les $F$ colossaux
restent interprétables comme ordres de grandeur)~; normalité des résidus
(TCL)~; \textbf{homoscédasticité} vérifiée par les tests de
\textbf{Bartlett} et de \textbf{Levene}. En cas de variances inégales,
bascule possible vers l'ANOVA de Welch.
\end{conditions}

\subsection{Résultats}
% Chiffres réels : F ~ 47, p ~ 1e-85 (1 facteur) ; interaction p ~ 5e-69
L'ANOVA à un facteur rejette massivement l'égalité des moyennes
(\(F\approx 47\), \(p\approx 10^{-85}\)). Le test de Tukey identifie
18~paires significatives sur 45, \texttt{rsi\_strict} se détachant du
reste. L'interaction stratégie~$\times$~régime est très significative
(\(p\approx 5\times10^{-69}\))~: \emph{la valeur d'une stratégie dépend
du régime de marché}.

% =====================================================================
\section{Gagner dépend-il du contexte ?}
\label{sec:b1-etape3}

\subsection{Présentation théorique}
% [À DÉVELOPPER] Test du khi-deux d'indépendance sur tableau de
% contingence : effectifs attendus E_ij = (N_i. N_.j)/N ;
% chi2 = somme (N_ij - E_ij)^2 / E_ij ~ chi2 à (r-1)(c-1) ddl.
\[
  \chi^2 = \sum_{i,j}\frac{(N_{ij}-E_{ij})^2}{E_{ij}},
  \qquad E_{ij}=\frac{N_{i\cdot}N_{\cdot j}}{N},
  \qquad \chi^2 \sim \chi^2_{(r-1)(c-1)}.
\]

\begin{conditions}
Effectifs théoriques \(E_{ij}\ge 5\) (règle de Cochran). Sur tableaux
\(2\times2\) ou à faibles effectifs, recours possible à la correction de
\textbf{Yates}, au \textbf{test exact de Fisher} ou au \textbf{G-test}.
Indépendance des observations~: même réserve qu'à la section
précédente (chapitre~\ref{chap:juger}).
\end{conditions}

\subsection{Résultats}
% Chiffres réels de etape3_chi2.txt
\begin{table}[H]
\centering
\caption{Dépendance de \texttt{win} au contexte (khi-deux d'indépendance).}
\label{tab:b1-etape3}
\small
\begin{tabular}{@{}lrrc@{}}
\toprule
Croisement & $\chi^2$ & ddl & Verdict \\
\midrule
\texttt{win} $\times$ \texttt{regime\_entry} & $112{,}3$   & 2  & lien significatif \\
\texttt{win} $\times$ \texttt{sector}        & $16{,}6$    & 10 & indépendance ($p=0{,}08$) \\
\texttt{win} $\times$ \texttt{strategy}      & $6\,381{,}7$& 9  & lien très fort \\
\bottomrule
\end{tabular}
\end{table}

Le taux de gain dépend fortement de la \emph{stratégie}, faiblement du
\emph{régime}, et pas du \emph{secteur}.

\begin{limites}
À très grand \(n\), le \(\chi^2\) devient significatif même pour des
liens négligeables~; la \emph{significativité} ne dit rien de la
\emph{force} du lien. % justifie la prudence d'interprétation
\end{limites}

% =====================================================================
\section{Le contexte d'entrée en carte --- ACP}
\label{sec:b1-etape4}

\subsection{Présentation théorique}
% [À DÉVELOPPER] ACP normée (Périnel) : centrage-réduction,
% diagonalisation de la matrice de corrélation R, valeurs propres
% lambda_k = inertie, cercle des corrélations, cos^2, CTR.
% Critères Kaiser / Cattell pour le nombre d'axes. Variables
% illustratives (régime, secteur, win) projetées a posteriori (V.test).
On réalise une \textbf{ACP normée}~\cite{saporta2006,lebart2006} sur les
variables quantitatives de contexte (\texttt{vol\_entry},
\texttt{rsi\_entry}, \texttt{dist\_ma200},
\texttt{holding\_days})~; le nombre d'axes est justifié par les critères
de \textbf{Kaiser} et de \textbf{Cattell}. Les variables qualitatives
(\texttt{regime}, \texttt{secteur}, \texttt{win}) sont projetées en
\textbf{illustratives} (valeur-test).

\begin{conditions}
Variables quantitatives~; standardisation obligatoire (unités
hétérogènes)~; l'ACP ne capture que des relations \emph{linéaires}.
\end{conditions}

\subsection{Résultats}
% [À COMPLÉTER avec etape4_acp.txt / etape4_acp.png]
\emph{[Insérer la figure \texttt{etape4\_acp.png}, l'inertie des axes et
la lecture du plan factoriel.]}

% =====================================================================
\section{AFC et ACM}
\label{sec:b1-etape56}

\subsection{Présentation théorique}
% [À DÉVELOPPER] AFC sur tableau de contingence strategy x tranche de
% résultat (profils lignes/colonnes, distance du chi-deux, résidus de
% Haberman). ACM sur plusieurs variables qualitatives (correction de
% Benzécri, rapport de corrélation eta^2).
L'\textbf{AFC} analyse le tableau de contingence croisant stratégie et
tranche de résultat~; l'\textbf{ACM} étend l'analyse à plusieurs
variables qualitatives simultanément.

\subsection{Résultats}
% [À COMPLÉTER avec etape5_afc.* et etape6_acm.*]
\emph{[Insérer les plans factoriels AFC et ACM et leur interprétation.]}

% =====================================================================
\section{Synthèse du chapitre}
\label{sec:contexte-synthese}
% [À RÉDIGER — 4-6 lignes] Idées :
%  - Les stratégies diffèrent fortement entre elles et leur valeur
%    dépend du régime de marché (interaction) ; le secteur ne joue pas.
%  - MAIS aucun de ces contrastes ne remet en cause le verdict du
%    chap. 3 : différer du voisin n'est pas battre le hasard.
%  - Ces cartes du contexte préparent la perspective « edge
%    conditionnel » (chap. 5).